\subsection{Dataset Description}

\paragraph{Long-form hallucination corpora.}
\begin{enumerate}
\item RAGTruth \citep{niu2024ragtruth} (License: MIT) is a
retrieval-augmented generation corpus with span-level hallucination
annotations across QA, summarization, and data-to-text prompts.
\item FactScore \citep{min2023} (License: MIT) is a
biography-generation corpus paired with atomic-fact verification
against Wikipedia.
\item LongFact \citep{wei2024} (License: Apache-2.0) is a
long-form factuality corpus annotated with the SAFE verifier of the
same paper.
\end{enumerate}

\paragraph{LM-Polygraph datasets.} We follow
\citet{vashurin2025lmpolygraph} for subset selection, prompt
formatting, and few-shot sources on the short-form benchmark
(TriviaQA, CoQA, MMLU, GSM8k, WMT14 Fr-En, WMT19 De-En, XSUM).

\paragraph{Splits.} For each corpus we follow the original protocol;
exact split sizes and prompt templates accompany the released
artefacts.

\paragraph{LLM-labeler accuracy on FactScore and LongFact.} The
auxiliary span labels $\tilde{y}_m$ that Stage 3 trains on are
produced by a claim-level verifier rather than human annotators.
We use Gemma 4 31B as the verifier on FactScore \citep{min2023} and
LongFact \citep{wei2024}, and we measure how closely its verdicts
match the human gold on a held-out subset of each corpus
(Table~\ref{tab:labeler}). The labeler's noise floor sets a ceiling
on what Stage 3 can hope to learn from the auxiliary signal.

\begin{table}[t]
\centering
\setlength{\tabcolsep}{4pt}
\renewcommand{\arraystretch}{1.05}
\scriptsize
\adjustbox{max width=\linewidth}{%
\begin{tabular}{l ccccc r}
\toprule
Dataset    & Accuracy & Precision & Recall & F1 & Cohen's $\kappa$ & $n$ (claims) \\
\midrule
FactScore  & ---      & ---       & ---    & ---& ---              & --- \\
LongFact   & ---      & ---       & ---    & ---& ---              & --- \\
\bottomrule
\end{tabular}%
}
\caption{\textbf{LLM-labeler agreement with human annotations} on
held-out subsets of FactScore \citep{min2023} and LongFact
\citep{wei2024}. Accuracy is the share of human-matching verdicts;
Precision/Recall/F1 use ``unsupported'' as the positive class
(Stage~3 supervision polarity); Cohen's $\kappa$ is chance-corrected
agreement. Protocol in Appendix~\ref{app:extended}. ``---'' = not
yet run.}
\label{tab:labeler}
\end{table}

